\documentclass{article}

\usepackage{fullpage}
\usepackage{amsmath}
\usepackage{amssymb}
\usepackage{fancyhdr}

\usepackage[colorlinks,linkcolor=dblue,filecolor=black,citecolor=dblue,urlcolor=dblue]{hyperref}
\usepackage{etex}
\bibliographystyle{alpha}
\usepackage[english]{babel}
\usepackage{amsfonts,amssymb,amsmath,sectsty,url}
\usepackage{mathrsfs}

\usepackage{tikz}
\usepackage{color}
%\usepackage{mathtools}
\usepackage[all]{xy}
\usepackage{ctable}
% \usetikzlibrary{calc}
\usepackage{tabularx,multirow}
\usepackage{graphicx}
\usepackage{ifpdf}
\usepackage{float}
\usepackage{multirow}
\usepackage{multicol,color}
\usepackage{tablefootnote}
\usepackage{pgfplots}

\usepackage{stmaryrd}
\usepackage{booktabs}
\usepackage{multirow}
\usepackage{import}
\usepackage{rotating}

% For tables
\usepackage{adjustbox}
\usepackage{booktabs}
\usepackage{multirow}
\usepackage{tabularx}

% For better precision math computations and formatting in tables
\usepackage{tikz,fp}
\usepackage[nottoc]{tocbibind}

\usepackage{listings}% http://ctan.org/pkg/listings
\lstset{
  basicstyle=\ttfamily,
  mathescape
}

% COMMANDS

\renewcommand{\emptyset}{\varnothing}

% Comments
\newcommand{\commentD}[1]{\textcolor[rgb]{0,1,0}{{\sf Drago\c{s}:} {\sl{#1}}}}
\newcommand{\commentM}[1]{\textcolor[rgb]{0,0.75,1}{{\sf Morten:} {\sl{#1}}}}
% Colors

\definecolor{dblue}{rgb}{0.00, 0.50, 0.90}
\definecolor{lblue}{rgb}{0.70, 0.80, 1.00}
\definecolor{lpink}{rgb}{0.90, 0.70, 1.00}
\definecolor{lgreen}{rgb}{0.80, 0.95, 0.75}
\definecolor{lred}{rgb}{0.99, 0.50, 0.55}
\definecolor{lyellow}{rgb}{1.00, 0.95, 0.75}
\definecolor{llgrey}{rgb}{0.95, 0.95, 0.95}
\definecolor{salmon}{rgb}{0.99, 0.90, 0.90}

% Dataflow stuff
\newcommand{\shape}{\mathsf{s}}

% UC security things

\newcommand{\Abort}{\mathsf{Abort}}
\newcommand{\Checkk}{\mathsf{Check}}
\newcommand{\Commit}{\mathsf{Commit}}
\newcommand{\Fcomm}{\mathcal F_\Commit}
\newcommand{\Fail}{\mathsf{Fail}}
\newcommand{\Keys}{\mathsf{Keys}}
\newcommand{\MACCheck}{\mathsf{MACCheck}}
\newcommand{\Open}{\mathsf{Open}}
\newcommand{\OK}{\mathsf{OK}}
\newcommand{\Pcomm}{\Pi_\Commit}
\newcommand{\Scomm}{\mathcal S_\Commit}
\newcommand{\Triple}{\mathsf{Triple}}
\newcommand{\Val}{\mathsf{Val}}

\newcommand{\VK}{\mathsf{vk}}
\newcommand{\DSB}{\mathsf{DSB}}
\newcommand{\DSA}{\mathsf{DSA}}
\newcommand{\OP}{\mathsf{op}}
\newcommand{\EK}{\mathsf{ek}}
\newcommand{\DK}{\mathsf{dk}}
\newcommand{\SK}{\mathsf{sk}}
\newcommand{\HASH}{\mathsf{hash}}
\newcommand{\SIGN}{\mathsf{Sign}}
\newcommand{\VER}{\mathsf{Ver}}
\newcommand{\TAG}{\mathsf{Tag}}

% fixed point arithmetic symbols
\newcommand{\Decode}{\mathsf{Decode}}


\newcommand{\inc}[1]{\mathsf{inc}({#1})}

% MPC things
\newcommand{\Pa}{\mathcal{P}}
\newcommand{\Engine}{\mathsf{Eng}}
\newcommand{\Ext}{\mathsf{Ext}}
\newcommand{\authshare}[1]{\llbracket {#1} \rrbracket} % a finite field
\newcommand{\auth}[1]{\authshare{#1}} % a finite field

% ZK things
\newcommand{\prov}{\mathcal{P}}
\newcommand{\ver}{\mathcal{V}}

% CS things
\newcommand{\Oh}{ \mathcal{O} }
\newcommand{\QuasiOh}{ \widetilde{ \mathcal{O} } }

\newcommand{\B}         {\mathrm{B}}
\newcommand{\KB}        {\mathrm{KB}}
\newcommand{\MB}        {\mathrm{MB}}
\newcommand{\GB}        {\mathrm{GB}}

\newcommand{\ms}        {\mathrm{ms}}
\newcommand{\mus}       {\mathrm{{\mu}s}}



% Sets, fields, rings, groups, etc...
\newcommand{\C}{\mathbb{C}} % complex nr
\newcommand{\Z}{\mathbb{Z}} % the integers
\newcommand{\Q}{\mathbb{Q}} % the rationals
\newcommand{\N}{\mathbb{N}} % the non-negative integers
\newcommand{\F}{\mathbb{F}} % a finite field
\newcommand{\Fp}{\mathbb{F}_p}
\newcommand{\Fn}{\mathbb{F}^n_{2}}
\newcommand{\Fq}{\mathbb{F}_q}
\newcommand{\range}[1]{[1\dots {#1} ]} % a finite field

\newcommand{\sample}{\stackrel{\small{\$}}{\leftarrow}}

\newcommand{\RingSample}{\mathsf{Sample}}
\newcommand{\SampleKey}{\mathsf{SampleKey}}
\newcommand{\dkey}{\textsf{dk}}
\newcommand{\GenNonce}{\textsf{GenNonce}()}

\newcommand{\M}{\mathcal{M}} % message space
\newcommand{\T}{\mathcal{T}} % tag space
\newcommand{\R}{\mathcal{R}} % randomness space

\newcommand{\tensor}{\otimes}
\newcommand{\unit}{\mathbf{u}}
\newcommand{\transpose}[1]{#1^{\intercal}}
\newcommand{\identitymatrix}{\mathbb{I}}



% Probability and DP
\newcommand{\Avg}{\textup{E}} % Average of random variables

\newcommand{\seed}{\mathsf{seed}} % Average of random variables
\newcommand{\derives}{\mathsf{derive\_seed}} % Average of random variables


\newcommand{\D}{\mathcal{D}} % a distribution
\newcommand{\U}{\mathcal{U}} % the uniform distribution

\newcommand{\uniform}[1]            {\mathcal{U}(#1)}
\newcommand{\binomial}[1]           {\mathcal{B}(#1)}
\newcommand{\gaussian}[1]           {\mathcal{N}(#1)}
\newcommand{\dgaussian}[1]     {\mathcal{DG}(#1)}
\newcommand{\hwt}[1]     {\mathcal{HWT}(#1)}
\newcommand{\zo}[1]     {\mathcal{ZO}(#1)}
\newcommand{\laplacian}[1]          {\mathrm{Lap}(#1)}
\newcommand{\distribution}          {\mathcal{D}}
\newcommand{\leftovernoise}     {\mathcal{E}}

\newcommand{\pr}[1]                 {\Pr \left[ #1 \right]}
\newcommand{\floor}[1]                  {\lfloor#1\rfloor}
\newcommand{\ceil}[1]                   {\left\lceil#1\right\rceil}



% Labels

\newcommand{\pk}{\mathsf{pk}}
\newcommand{\sk}{\mathsf{sk}}
\newcommand{\sfT}{\mathsf{T}}

\newcommand{\Sim}{\mathcal{S}}
\newcommand{\Simu}[1]{\mathcal{S}_{\mathsf{#1}}}



% Operators

\newcommand{\asn}{\gets} % to be read as ``computed as'' or ``taken according to the distribution...'' or ``taken uniformly from...''

\newcommand{\ParamGen}{\mathsf{ParamGen}}
\newcommand{\KeyGen}{\mathsf{KeyGen}}
\newcommand{\Enc}{\mathsf{Enc}}
\newcommand{\Dec}{\mathsf{Dec}}
\newcommand{\Add}{\mathsf{Add}}
\newcommand{\ek}{ek}
\newcommand{\dk}{dk}
\newcommand{\plaintext}{m}
\newcommand{\ciphertext}{e}
\newcommand{\plaintextspace}{\mathcal{P}}
\newcommand{\ciphertextspace}{\mathcal{E}}
\newcommand{\INDCPA}{\textrm{IND-CPA}}

\newcommand{\Share}{\mathsf{Share}}
\newcommand{\Shamir}{\mathsf{Shamir}}
\newcommand{\RS}{\mathsf{RS}}
\newcommand{\Rec}{\mathsf{Recon}}

\newcommand{\share}[1]{\langle {#1} \rangle} % a finite field
\newcommand{\shareR}[1]{{\langle {#1} \rangle}_{2^k}} % a ring
\newcommand{\shareB}[1]{{\langle {#1} \rangle}_2} % a binary field
\newcommand{\secret}{m}

\newcommand{\MAC}{\mathsf{MAC}}
\newcommand{\Encode}{\mathsf{Encode}}

\newcommand{\PRG}{\mathsf{PRG}}

\newcommand{\Adv}{\mathsf{Adv}}
\newcommand{\Gen}{\mathsf{Gen}}



\newcommand{\ineff}{\mathsf{ineff}}
\newcommand{\superpoly}{\mathsf{superpoly}}
\newcommand{\poly}{\mathsf{poly}}
\newcommand{\negl}{\mathsf{negl}}
\newcommand{\const}{\mathsf{const}}
\newcommand{\polylog}{\mathsf{polylog}}

\newcommand{\normtwo}[1]{\|{#1}\|_2}
\newcommand{\normI}[1]{\|{#1}\|_\infty}

\newcommand{\Bclean}{B_\mathsf{clean}}
\newcommand{\Bplain}{B_\mathsf{plain}}
\newcommand{\Brand}{B_\mathsf{rand}}

\newcommand{\secp}{\mathsf{sec}} % statistical security parameter
\newcommand{\ssec}{\secp} % use this for stat sec param
\newcommand{\csec}{k} % comp sec param

\newcommand{\Func}{\mathcal{F}}
\newcommand{\Funct}[1]{\Func_\mathsf{#1}}
\newcommand{\Proto}[1]{\Pi_\mathsf{#1}}
\newcommand{\nPi}{\mathrm{\Pi}}
\newcommand{\nPhi}{\mathrm{\Phi}}
\newcommand{\Bra}[1]{\llbracket #1 \rrbracket}
\newcommand{\Fauth}{\Func_{\Bra{\cdot}}}
\newcommand{\Pauth}{\Proto{\Bra{\cdot}}}
\newcommand{\Sauth}{\Sim_{\Bra{\cdot}}}
\newcommand{\PMACCheck}{\Proto{\MACCheck}}
\newcommand{\Ftriple}{\Funct{\Triple}}
\newcommand{\Ptriple}{\Proto{\Triple}}
\newcommand{\Striple}{\Sim_{\Triple}}
\newcommand{\Fmult}{\Funct{Mult}}
\newcommand{\Pmult}{\Proto{2Mult}}
\newcommand{\Smult}{\Sim_\mathsf{2Mult}}
\newcommand{\Ffmult}{\Funct{2FMult}}
\newcommand{\Pfmult}{\Proto{2FMult}}
\newcommand{\Sfmult}{\Sim_\mathsf{2FMult}}
\newcommand{\Frand}{\Funct{Rand}}
\newcommand{\Fpok}{\Funct{ZKPoK}^S}
\newcommand{\Fgpok}{\Funct{gZKPoK}^S}
\newcommand{\Pgpok}{\Proto{gZKPoK}^S}
\newcommand{\Sgpok}{\Simu{gZKPoK}^S}
\newcommand{\Gcpa}{\mathcal{G}_\mathsf{cpa+}}

\newcommand{\vl}{\mathbf{l}}
\newcommand{\vL}{\mathbf{L}}
\newcommand{\vz}{\mathbf{z}}
\newcommand{\vx}{\mathbf{x}}
\newcommand{\vm}{\mathbf{m}}
\newcommand{\vM}{\mathbf{M}}
\newcommand{\vy}{\mathbf{y}}
\newcommand{\vu}{\mathbf{u}}
\newcommand{\vk}{\mathbf{k}}
\newcommand{\vq}{\mathbf{q}}
\newcommand{\vt}{\mathbf{t}}
\newcommand{\vs}{\mathbf{s}}
\newcommand{\vr}{\mathbf{r}}
\newcommand{\vv}{\mathbf{v}}
\newcommand{\va}{\mathbf{a}}
\newcommand{\vh}{\mathbf{h}}
\newcommand{\vb}{\mathbf{b}}
\newcommand{\vc}{\mathbf{c}}
\newcommand{\vd}{\mathbf{d}}
\newcommand{\vdp}{\mathbf{d^{'(i)}}}
\newcommand{\ve}{\mathbf{e}}
\newcommand{\vw}{\mathbf{w}}
\newcommand{\vX}{\mathbf{X}}
\newcommand{\vC}{\mathbf{C}}
\newcommand{\vn}{\mathbf{n}}
\newcommand{\vf}{\mathbf{f}}
\newcommand{\vg}{\mathbf{g}}

\newcommand{\Input}{\mathsf{Input}}
% \newcommand{\Triple}{\mathsf{Triple}}
\newcommand{\Lincomb}{\mathsf{LinComb}}
\newcommand{\id}{\mathsf{id}}

\newcommand{\vDelta}{\boldsymbol{\Delta}}

\newcommand{\valpha}{\boldsymbol\alpha}
\newcommand{\vbeta}{\boldsymbol\beta}
\newcommand{\vchi}{\boldsymbol\chi}
\newcommand{\vdelta}{\boldsymbol\delta}
\newcommand{\vepsilon}{\boldsymbol\epsilon}
\newcommand{\veta}{\boldsymbol\eta}
\newcommand{\vgamma}{\boldsymbol\gamma}
\newcommand{\vlambda}{\boldsymbol\lambda}
\newcommand{\vmu}{\boldsymbol\mu}
\newcommand{\vnu}{\boldsymbol\nu}
\newcommand{\vphi}{\boldsymbol\phi}
\newcommand{\vpsi}{\boldsymbol\psi}
\newcommand{\vrho}{\boldsymbol\rho}
\newcommand{\vsigma}{\boldsymbol\sigma}
\newcommand{\vtau}{\boldsymbol\tau}
\newcommand{\vtheta}{\boldsymbol\theta}
\newcommand{\vxi}{\boldsymbol\xi}
\newcommand{\vzeta}{\boldsymbol\zeta}

\newcommand{\ui}{^{(i)}}
\newcommand{\upi}{^{'(i)}}
\newcommand{\uti}{^{(i){\intercal}}}

\newcommand{\uj}{^{(j)}}
\newcommand{\uij}{^{(ij)}}
\newcommand{\uji}{^{(ji)}}

\newcommand{\mask}[1]{\ifnum\ANON=0#1\fi}
\newcommand{\anon}[1]{\ifnum\ANON=1#1\fi}
\newcommand{\maskanon}[2]{\ifnum\ANON=1#1\else#2\fi}




% Environments
\newenvironment{boxfig}[2]{% {#1}{#2} = {Caption}{label}
\begin{figure}
  \newcommand{\FigCaption}{#1}
  \newcommand{\FigLabel}{#2}
  \vspace{-\medskipamount}
  \begin{center}
    \begin{small}
      \begin{tabular}{@{}|@{~~}l@{~~}|@{}}
        \hline
        \rule[-1.5ex]{0pt}{1ex}
        \begin{minipage}[b]{.96\linewidth}
          \vspace{1ex}
          \smallskip
          }{%
        \end{minipage}\\
        \hline
      \end{tabular}
    \end{small}
    \vspace{-1\bigskipamount}
    \caption{\small \FigCaption}
    \label{\FigLabel}
%  \vspace{-0.3cm}
  \end{center}
\end{figure}
}
\newenvironment{Boxfig}[3]{% {Caption}{Label}{Title}
  \begin{boxfig}{#1}{#2}
    \begin{center}
      \textbf{#3}
    \end{center}
  }{%
    \end{boxfig}
  }



% LATEX

% this crossreferences description items
\makeatletter
\let\orgdescriptionlabel\descriptionlabel
\renewcommand*{\descriptionlabel}[1]{%
  \let\orglabel\label
  \let\label\@gobble
  \phantomsection
  \edef\@currentlabel{#1}%
  %\edef\@currentlabelname{#1}%
  \let\label\orglabel
  \orgdescriptionlabel{#1}%
}
\makeatother




% a temporary counter
\newcounter{tempcount}

%%% Local Variables:
%%% mode: latex
%%% TeX-master: "main"
%%% End:

%\usepackage[color=oxygenorange]{todonotes}
%\newcommand{\crec}[1]{\todo[inline]{\textbf{crec:} #1}\xspace}

% \usepackage{draftwatermark}
% \SetWatermarkLightness{0.9}
% \SetWatermarkText{DRAFT \\ \today}
% \SetWatermarkScale{2}


\usetikzlibrary{matrix,chains,positioning,decorations.pathreplacing,arrows,calc}

\newcommand\tikzmark[2]{%
\tikz[remember picture,baseline] \node[right, outer sep=0pt, inner sep=0pt] (#1){\phantom{#2}};%
}

\newcommand\link[3]{%
\begin{tikzpicture}[remember picture, overlay, >=stealth, shift={(0,0)}]
  \draw[->] (#1) -- node[auto,] {#3} (#2);
\end{tikzpicture}%
}

% \pagestyle{fancy}
% \fancyhf{}
% \setlength{\headheight}{22pt}
% \setlength{\headsep}{0.2in}
% \renewcommand{\headrulewidth}{0.4pt}
% \renewcommand{\footrulewidth}{0.4pt}% default is 0pt
% \cfoot{\thepage}
% \lhead{\includegraphics[width=3cm]{cape_logo.png}}
% \rhead{Proprietary and Confidential}

\newcommand\todo[1]{\textcolor{red}{#1}}

%
\begin{document}
%
\newcommand{\mainsection}[1]{\newpage \section{#1}}
\newcommand{\msubsection}[1]{\newpage \subsection{#1}}
\newcommand{\msubsubsection}[1]{\subsubsection{#1}}

\title{$\mathsf{Moose}$ documentation}
% \author{Morten Dahl, Yann Dupis, Jason Mancuso, Dragos Rotaru}
% \institute{Cape Privacy}


\maketitle
\tableofcontents

\thispagestyle{fancy}


\section{Introduction}

This paper details the cryptographic techniques and protocols used by Cape
Privacy's encrypted learning solution. We expect the reader to have some
familiarity with secure multi-party computation, linear algebra, and ring
arithmetic; as a good starting point we recommend \cite{evans2017pragmatic}. We
also omit detailed explanations of machine learning and the models we support,
instead referring the interested reader to~\cite{mml}.

The main purpose of this paper is to describe some of the design decisions that went into Moose
and the current MPC protocols that are implemented in it.

% secure linear
% regression. Future versions of the paper will expand upon this with additional
% models and protocols.

% uncomment this for linear regression
% 
\subsection{Linear regression}
\label{sec:linreg}

We focus on performing linear regression between two distinct data owners, one providing the explanatory variable $X$ and another providing some target (or response) variable $y$. For simplicity, we assume $X$ is an $n \times p$ matrix with $n > p$ and $y$ a column vector, indexed row-wise by $X_i$ and $y_i$, respectively.

Solving the classical linear regression problem equates to finding the minimizer of the mean squared error objective
\begin{align*}
\optCoeff = \min_{\coeff}\instanceAvg (X_i \cdot w - y_i)^2
\end{align*}
Using the method of maximum likelihood estimation~\cite{mml}, we can reformulate this optimization problem into the product
\begin{align*}
\optCoeff &= (X^TX)^{-1}X^Ty
\\
&= Z \cdot y
\end{align*}
Note that evaluating $Z = (X^TX)^{-1}X^T$ can be performed as local preprocessing on plaintext data by the provider of $X$ (since the righthand term does not involve the target variable $y$). Thus, solving this regression problem securely reduces to a single secret-shared matrix-vector product $\coeff = Z \cdot y$.

Additionally, we may want to compute certain metrics to assess the regression
outcome. These metrics must also be evaluated on secret-shared data. Concretely,
we present cryptographic protocols to support mean squared error
($\mathsf{MSE}$), mean absolute percentage error
($\mathsf{MAPE}$), and R-squared ($R^2$) metrics. Denote the target predictions $\ypred = X \cdot
\coeff$ and the expected target value as $\bar{y} = \instanceAvg y_i$. An
argument similar to the above reduces the secure computation of $R^2$ to
securely computing the residual sum of squares ($\mathsf{RSS}$):
\begin{align*}
R^2 &= 1 - \frac{\sum(y_i - \ypred_i)^2}{\sum(y_i - \bar{y})^2}
\\
&= 1 - \frac{RSS(\ypred, y)}{SS(\ypred, y)}
\end{align*}
since only the $RSS$ term depends on data from both providers (i.e. both $X$ and $y$).

In summary, we are ultimately concerned with providing protocols for securely evaluating the following values:

\begin{align*}
\coeff &= Z \cdot y
\\
\ypred &= X \cdot \coeff
\\
\mathit{MSE}(\ypred, y) &= \instanceAvg (\ypred - y_i)^2
\\
\mathit{RSS}(\ypred, y) &= \sum_n (\ypred - y_i)^2
\\
\mathsf{MAPE}(\ypred, y) &= \frac{1}{n} \sum_n \Big| \frac{\ypred - y_i}{y_i} \Big|
\end{align*}

%In order to compute the MAPE (mean absolute percentage error) metric which boils down to compute the absolute value of a secret $\share{|\vx|}$.

% \definecolor{codegreen}{rgb}{0,0.6,0}
% \definecolor{codegray}{rgb}{0.5,0.5,0.5}
% \definecolor{codepurple}{rgb}{0.58,0,0.82}
% \definecolor{backcolour}{rgb}{0.95,0.95,0.92}

% \lstdefinestyle{mystyle}{
%     backgroundcolor=\color{backcolour},
%     commentstyle=\color{codegreen},
%     keywordstyle=\color{magenta},
%     numberstyle=\tiny\color{codegray},
%     stringstyle=\color{codepurple},
%     basicstyle=\ttfamily\footnotesize,
%     breakatwhitespace=false,
%     breaklines=true,
%     captionpos=b,
%     keepspaces=true,
%     numbers=left,
%     numbersep=5pt,
%     showspaces=false,
%     tabsize=2
% }

% \lstset{style=mystyle}

% \begin{lstlisting}[language=Python]
% @edsl.computation
% def lin_reg():

%     with x_owner:
%         X = edsl.load("X")
%         bias = edsl.ones(edsl.slice(edsl.shape(X), begin=0, end=1))
%         reshaped_bias = edsl.expand_dims(bias, 1)
%         X_b = edsl.concatenate([reshaped_bias, X], axis=1)
%         A = edsl.inverse(edsl.dot(edsl.transpose(X_b), X_b))
%         B = edsl.dot(A, edsl.transpose(X_b))

%     with y_owner:
%         y_true = edsl.load("y")
%         totals_ss = ss_tot(y_true)

%     with replicated:
%         w = edsl.dot(B, y_true)
%         y_pred = edsl.dot(X_b, w)
%         mse_result = mse(y_pred, y_true)
%         residuals_ss = ss_res(y_pred, y_true)

%     with x_owner:
%         rsquared_result = r_squared(residuals_ss, totals_ss)
%         w = edsl.identity(w)
%         mse = edsl.identity(mse)
%         residuals_ss = edsl.identity(residuals_ss)
% \end{lstlisting}

Finally, these operations must be implemented for fixed-point numbers as described in Section~\ref{sec:fixed-point}.


\subsection{Fixed point arithmetic}
\label{subsec:fixed-point}

We define $\verb|fixed|(k, f)$ as the set of rational numbers $\{x \in \Q : x
= \bar{x} \cdot 2^{-f}, \bar{x} \in \Z_{\share{k}}\}$. Here $\bar{x} \in
Z_{\share{k}}$ denotes that $\bar{x}$ is at most a $k$ bit integer. A fixed
point number is represented in memory as $x \cdot 2^{f} = \bar{x} \in
\Z_{\share{k}}$. There are a few generic protocols to encode, decode or add
and multiply two fixed point numbers but we focus here on truncation while
the interested reader can check more details in
Appendix~\ref{app:fixed-point}.

\subsection{Truncation}
Given secret shared $\share{a}$ the goal is to compute $\share{a \bmod 2^m}$
where $m$ is a public constant. This is also denoted as
$\mathsf{trunc}(\share{a}, m)$. There are many ways to truncate a secret that
can be roughly classified into two lines of thought: a) probabilistic
truncation where there is a rounding error with probability $p$ and b) deterministic
truncation where the output is always exact.

Current state of the art techniques in truncating a secret use
dabits/edabits as described by Escudero et al. in \cite{C:EGKRS20} and
various share conversions. In \cite{C:EGKRS20} there are several techniques
to achieve both probabilistic and deterministic truncation for field and ring
case. To avoid edabit generation required for generating random bounded integers
by the truncation protocols in \cite{C:EGKRS20} we use the protocols from
\cite{PoPETS:DalEscKel20}. We describe their probabilistic protocols for replicated secret
sharing below in Fig~\ref{fig:truncpr-rep-ss}.
Note that although the protocol in Fig~\ref{fig:truncpr-rep-ss} looks
very efficient, the edabit variant from \cite{C:EGKRS20} requires
about $\approx 35$ times less communication. \commentD{Need to re-check this, I think it's not that bad}



\begin{Boxfig}{Truncation protocol for semi-honest
RSS.}{fig:truncpr-rep-ss}{Protocol $\Proto{\mathsf{TruncPr}}$}
On input
$\mathsf{TruncPrSp}(\share{\vx}, m)$ from all parties where $\vx$ is a sharing
of a tensor over $\Z_{2^k}^{\mathsf{s}}$: \\
$P_3$ does the following:
  \begin{enumerate}
    \item Sample key $k_3 \sample \mathsf{SampleKey}()$.
    \item For $i \in [k]$ call $\seed(i) \asn \derives(k_3, i)$
    \item For $i \in [k]$ sample locally random bits $\mathbf{b}_i \in \Z_{2^k}^{\mathsf{s}}$.
    This is done by calling
    generate $\mathbf{b}_i \sample \mathsf{Sample}(\seed(i), \mathsf{s}, [0,1])$.
    \item Compute
    $\mathbf{r}_{\mathsf{bot}} = \sum_{i=0}^{k-1}2^i \cdot \vb(i)$,
    $\mathbf{r}_{\mathsf{msb}} \asn \vb(k-1)$ and
    $\mathbf{r}_{\mathsf{top}} = \sum_{i=m}^{k-2} \vb(i) \cdot 2^{i-m}$. Send
    the corresponding shares to $P_1$ and $P_2$.
    \item For each of $\vr_{\mathsf{bot}}, \vr_{\mathsf{msb}} \vr_{\mathsf{top}}$
    generate 2-out-of-2 additive shares by calling $\Proto{\mathsf{AdditiveShare}}(\cdot)$
    on each entry and get $
    \vr\ui_{\mathsf{bot}}, \vr\ui_{\mathsf{msb}}, \vr\ui_{\mathsf{top}}$ for $i \in [0,1]$. Send the corresponding $\vr\ui$ to party $P_i$.
 \end{enumerate}
$P_1$ together with $P_2$ do the following:

\begin{enumerate}
   \item Convert $\share{\vx}$ into a $2$-out-of-$2$ sharing.
   \begin{enumerate}
      \item On $P_1$ set $\vx'_1 \asn \vx_1 + \vx_2$.
      \item On $P_2$ set $\vx'_2 \asn \vx_2$.
   \end{enumerate}

   \item $P_1$ and $P_2$ then execute the protocol $\Proto{\mathsf{TruncPR}}(\share{\vx'}, m)$
   using the $2$-out-of-$2$ sharing of $\vx$. Store the output as $\vy$.
\end{enumerate}

 Now parties convert the $2$-out-of-$2$ sharing $\vy$ to a $2$-out-of-$3$
 sharing. All parties output $\Proto{(2,2)\rightarrow(2,3)}(\vy)$.

%  \begin{enumerate}
%    \item Each $P_i$ for $i \in [1, 2]$ sets the output of $\mathsf{TruncPr}$
%    as $y'_i$ and $\hat{y}_i$ as the sharing received from $P_3$. Send $y'_i -
%    \hat{y}_i$ to $P_{3-i}$. The received values are denoted by $\tilde{y}_i$.
%    \item $P_1$ sets the output share as $(y_1, y'_1 - \hat{y}_1 + \tilde{y}_1)$
%    \item $P_2$ sets the output share as $(y'_2 - \hat{y}_2 + \widehat{y}_2, y_3)$.
%  \end{enumerate}
\end{Boxfig}

\begin{Boxfig}{Generic truncation protocol.}{fig:truncpr-generic}
  {Protocol $\Proto{\mathsf{2PCTruncPr}}(\share{\vx}, m)$}
  Note this is a two party protocol, meaning only $P_1$ and $P_2$ are involved.
  Assume $\vx \in \Z_{2^k}^{\shape}$ has $\mathsf{MSB}(\vx) = 0$. If $\vx$ is negative then
  we make it positive by adding $2^{k-1}$ to it and assume that $\mathsf{abs}(\vx) < 2^{k-1}$.
  \\
  On $P_1$ and $P_2$:
  \begin{enumerate}
    \item Compute $\vc = \Open(\share{\vx} + \share{\vr_{\mathsf{bot}}})$.
    \item $\vc' \asn (\vc/2^m) \bmod 2^{k-m-1}$. This can be done using one $\mathsf{ring\_shl}$
    followed by one $\mathsf{ring\_shr}$ operation.
    \item Compute $\share{\vb} \asn \share{\vr_{\mathsf{msb}}} \oplus (\vc / 2^{k-1})$. Note here that $\oplus$
    in arithmetic circuits boils down to $a + b - 2a\cdot b$. Here one of the operands is public so
    the result can be computed using only local operations.
    \item Output $\vc' - \share{\vr_{\mathsf{top}}} + \share{\vb} \cdot 2^{k-m-1}$.
 \end{enumerate}

\end{Boxfig}


\begin{Boxfig}{Local conversion to $(2,2)$ additive shares}{fig:get-additive}
  {Protocol $\Proto{\mathsf{AdditiveShare}}(\vx)$}
  Roles: D, $P_1$, $P_2$ where $\vx.\mathsf{placement}$ is Host(D) \newline
  With Host(D):
  \begin{enumerate}
    \item $k \asn \mathsf{SampleKey}()$.
    \item $\seed \asn \derives(k, 0)$.
    \item $\vx^1 \asn \mathsf{SampleUniform}(\vx.\mathsf{shape}, \seed)$.
    \item $\vx^2 \asn \vx - \vx^1$.
    \item Send $\vx^i$ to $P_i$.
  \end{enumerate}
  With Host($P_1$):
  \begin{enumerate}
      \item Receive $\vx^1$ from Dealer.
  \end{enumerate}
  With Host($P_2$):
  \begin{enumerate}
      \item Receive $\vx^2$ from Dealer.
  \end{enumerate}
  Let $\vx = Additive(\vx^1, \vx^2)$.
\end{Boxfig}

\begin{Boxfig}{Opening protocol}{fig:open2-2}
  {Protocol $\Proto{\Open}^{(2,2)}(\share{\vx})$}
  On $P_1$:
  \begin{enumerate}
    \item Send $\vx_1$ to $P_2$.
  \end{enumerate}
   On $P_2$:
  \begin{enumerate}
    \item Send $\vx_2$ to $P_1$.
  \end{enumerate}
  On $P_1$
  \begin{enumerate}
    \item Receive $\vx_2$ from $P_2$. Output
    $\vx \asn \vx_1 + \vx_2$.
  \end{enumerate}
  On $P_2$
  \begin{enumerate}
    \item Receive $\vx_1$ from $P_1$. Output $\vx \asn \vx_1 + \vx_2$.
  \end{enumerate}
\end{Boxfig}


% \begin{Boxfig}{Opening protocol}{fig:open2-3}
%   {Protocol $\Proto{\mathsf{OpenTo}}^{(2,3)}(\share{\vx})$}
%   On $P_i$:
%   \begin{enumerate}
%     \item Send $\vx_1$ to $P_2$.
%   \end{enumerate}
%    On $P_2$:
%   \begin{enumerate}
%     \item Send $\vx_2$ to $P_1$.
%   \end{enumerate}
%   On $P_1$
%   \begin{enumerate}
%     \item Receive $\vx_2$ from $P_2$. Output
%     $\vx \asn \vx_1 + \vx_2$.
%   \end{enumerate}
%   On $P_2$
%   \begin{enumerate}
%     \item Receive $\vx_1$ from $P_1$. Output $\vx \asn \vx_1 + \vx_2$.
%   \end{enumerate}
% \end{Boxfig}







\section{Computational Model}

In this section we describe our computational model for protocols, which is loosely based on the data-flow paradigm~\cite{abadi16} and a simplified UC model~\cite{canetti-cohen-lindell15}. Concretely, all protocols are expressed as data-flow graphs where each node represents an operation to be performed by a specific party, and each edge presents a value ``flowing'' between operations.

The main motivation for using the data-flow paradigm is that it leads to a very natural concurrent execution model, where in the extreme we can see each node as being executed by a separate task (e.g. green thread or actor). We take full advantage of this in the runtime which is based on the asynchronous execution paradigm. The reason for basing our computational model on the UC model is that it is a well-known paradigm for ensuring security under concurrent composition.


\subsection{Sessions}

Every execution of a graph is performed under a unique session id $\sid$ used to identify values and ensure isolation when running protocols concurrently. As we shall see in more detail later, session ids are for instance used to non-interactively derive nonces and sample correlated randomness by the secret sharing schemes.

There are no requirements on session ids besides uniqueness, and it can safely be chosen by an untrusted coordinator, for instance by sampling a random string. However, for security reasons, each party maintains a list of previous session ids in which it has engaged, and refuses to re-run any computation using those ids. This prevents for instance replay and selective failure attacks.


\subsection{Sub-protocols}

We allow graphs to call sub-graphs similar to sub-routines. When doing so, the sub-graph is executed under a unique sub-session id $\sid'$ derived from $\sid$ and an activation key $\mathsf{ac\_key}$ statically related to the call site: 
$$
\sid' = \sid \| \mathsf{ac\_key}
$$
with $\|$ denoting string concatenation. In the presentation given in this paper, we make heavy use of calling sub-protocols, yet we typically (including in the current implementation) inline their sub-graphs for performance reasons. To do so securely, the activation key associated with the call site is embedded into the sub-graphs's operations during inlining.

As an example, $\mathsf{DeriveSeed}$ operations have a $\mathsf{sync\_key}$ attribute used to control uniqueness within the scope of the sub-graph, and during inlining this attribute is updated by prefixing it with the activation key, $\mathsf{sync\_key}' = \mathsf{ac\_key} \| \mathsf{sync\_key}$, so that the uniqueness is maintained in the combined scope.

This approach was chosen to ensure that we (in the future) can build a deterministic compilation process of computations, yet still support inlining.


\subsection{Communication}

Transmission of values between parties is done using $\mathsf{Send}$ and $\mathsf{Receive}$ nodes, each with a static rendezvous key $\mathsf{rdv\_key}$ attribute linking them together. Together with the session id, this allows us to uniquely identify all values by tagging them with $(\sid, \mathsf{rdv\_key})$, and is used to derive the nonce under which messages are encrypted as well as by the Cape infrastructure to keep track of messages (see Section~\ref{sec:infrastructure}).

The rendezvous key is also updated during inlining by prefixing it with the call site's activation key, ensuring that rendezvous keys likewise maintain their uniqueness: $\mathsf{rdv\_key}' = \mathsf{ac\_key} \| \mathsf{rdv\_key}$.

\import{replicated}{replicated.tex}
\section{Implementation}


\subsection{Runtime}

Operations on ring tensors are implemented using the \texttt{ndarray} Rust crate on
\texttt{Wrapping<u64>} and \texttt{Wrapping<u128>} scalars. Computations are executed
asynchronously via the \texttt{tokio} Rust crate by spawning a task for each operations.


\subsection{AES based PRG}
\label{sec:aes}

To efficiently sample a large batch of random bits, we use an AES-based PRG. We can produce $n \cdot 128$ random bits by calling $\mathsf{AES}_k(0), \cdots, \mathsf{AES}_k(n-1)$ with a $128$ bit key $k$.
Since AES works on $128$ bit blocks, the integer counters $0, \cdots, (n-1)$
are encoded using Little Endian format. In case one wants to retrieve a
number of $\ell$ random bits which is non-divisible by $128$, the last
$128 - \ell$ bits are discarded.

This method of generating
randomness is used in other well-known MPC libraries, such as MP-SPDZ
\cite{CCS:Keller20}, SCALE-MAMBA \cite{SCALE} or Swanky \cite{url-swanky}.
Security proofs for using fixed key AES in various MPC protocols were given
in \cite{SP:GKWY20}.

The PRNG wrapper we have built can be found in \textsf{rust/src/prng.rs} and
it is based on \cite{url-cargo-aes}.  We set the \textsf{AES} crate to
compile with special SSE3 CPU instructions whenever possible.
The PRNG key $k$ is generated using a call to \textsf{randombytes\_into}
from the Sodium Oxide rust crate \cite{url-sodiumoxide}.


% \import{key-distribution}{infra.tex}


\newpage
% ---- Bibliography ----
%
\bibliography{mybib,cryptobib/abbrev3,cryptobib/crypto}


\clearpage
\appendix


% \begin{Boxfig}{Sampling functionality}{fig:Sample}{Functionality
$\Func_{\mathsf{Sample}}$}
\textbf{Additive}: On input $(\mathsf{SampleAdd}, d, S, P_i)$ from party
$P_i$ the functionality samples a random tensor $\vx \sample S^d$. It also
generates two keys $k_1, k_2$ associated to $\vx$ such that when expanded using PRGs
they sum up to $\vx$, i.e. $\PRG(k_1) + \PRG(k_2) = \vx$.
\end{Boxfig}

% \import{key-distribution}{protocol-flow}
\import{replicated}{extended-protos}
\end{document}



